\section{语言模型基础：从自动补全说起}

\subsection{什么是语言模型}

\begin{importantbox}{语言模型的本质}
语言模型（Language Model）本质上是一个``高级自动补全''系统（fancy autocomplete）。给定一段文本前缀（stem），模型预测下一个最可能出现的词（token）。
\end{importantbox}

讲师从最直观的例子出发：给定 ``It's raining cats and \_\_\_\_''，几乎所有人都会预测下一个词是 ``dogs''。这正是语言模型要做的事情——根据上下文预测下一个词。

这个过程可以扩展为连续预测多个词。例如，给定 ``To be or not''，模型先预测 ``to''，将其拼接到上下文末尾，再预测 ``be''，从而恢复完整的短语 ``To be or not to be''。这种逐个 token 预测并反馈的方式被称为\textbf{自回归解码}（Autoregressive Decoding）。

\begin{knowledgebox}{自回归解码（Autoregressive Decoding）}
自回归解码是当前主流 LLM 生成文本的基本范式：
\begin{enumerate}
    \item 给定输入序列 $x_1, x_2, \dots, x_t$
    \item 模型预测下一个 token 的概率分布 $P(x_{t+1} \mid x_1, \dots, x_t)$
    \item 从分布中采样一个 token $x_{t+1}$
    \item 将 $x_{t+1}$ 拼接到序列末尾，重复步骤 2-3
    \item 直到生成特殊的终止 token（end-of-sentence）
\end{enumerate}
\end{knowledgebox}

\subsection{将问题嵌入``填空''结构}

语言模型之所以引发巨大兴趣，关键在于各种类型的任务都可以被巧妙地嵌入（embed）到这个``填空''问题结构中：

\begin{itemize}
    \item \textbf{数学问题}：``I have two apples and I eat one, I'm left with \_\_\_\_'' $\rightarrow$ ``one'' ——模型做了算术
    \item \textbf{类比推理}：``Paris is to France as Tokyo is to \_\_\_\_'' $\rightarrow$ ``Japan'' ——模型做了类比
    \item \textbf{事实查询}：``Pizza was invented in \_\_\_\_'' $\rightarrow$ ``Naples, Italy'' ——模型做了知识检索
\end{itemize}

讲师特别指出，类比推理曾长期困扰 NLP 研究者。在 Word2Vec embeddings 时代，研究者费尽心思也很难让模型解好类比题，而高中生却能在 SAT 考试中轻松应对。这一鲜明对比凸显了 LLM 带来的能力飞跃。

\subsection{构建一个统计语言模型}

讲师用 Bayesian（n-gram）语言模型来展示最基本的语言建模方法。这种方法早在 1980 年代就已提出，其核心思想是``fancy counting''（高级计数）。

以狄更斯的名句 ``It was the best of times, it was the worst of times'' 为语料库：

\begin{enumerate}
    \item \textbf{文本预处理}：转小写、去除标点、添加句首（start-of-sentence）和句尾（end-of-sentence）特殊 token
    \item \textbf{构建 n-gram 字典}：统计所有 unigram、bigram、trigram、4-gram 出现的次数
    \item \textbf{条件概率估计}：对于 stem ``it was the''，统计后续出现的词及频率
\end{enumerate}

给定 stem ``it was the''，语料中后续词的概率分布为：

\begin{table}[H]
\centering
\begin{tabular}{lcc}
\toprule
\textbf{下一个词} & \textbf{出现次数} & \textbf{概率} \\
\midrule
age       & 2 & 1/3 \\
best      & 1 & 1/6 \\
epoch     & 2 & 1/3 \\
worst     & 1 & 1/6 \\
\bottomrule
\end{tabular}
\caption{基于 n-gram 统计的条件概率分布}
\end{table}

通过从这个概率字典中随机采样，可以将该语言模型转化为一个\textbf{生成模型}。但生成结果往往会陷入``概率循环''（probability loop）：

\begin{quote}
\textit{``It was the best of times, it was the worst of times, it was the worst of times, it was the worst of times, it was the age of wisdom, it was the age of foolishness...''}
\end{quote}

\begin{warningbox}{概率循环与幻觉的直觉联系}
这个简单例子展示了一个重要现象：模型陷入了重复循环，不断输出 ``it was the worst of times''。这不是模型``悲观''，而是上下文窗口太小，导致模型不知道何时``跳出''重复模式。当你听到 LLM ``hallucinating''（幻觉）时，底层机制与此类似——模型在概率分布的奇异区域中生成文本，不知道该说什么，于是随机产生看似合理但实际错误的内容。
\end{warningbox}

\subsection{本章小结}

语言模型的核心任务是\textbf{下一词预测}（next-word prediction）。从最简单的 n-gram 统计模型到现代 Transformer LLM，这一基本范式没有改变。所有看似``智能''的表现——做数学、解类比、查事实——本质上都是通过将问题嵌入``填空''结构来实现的。n-gram 模型的局限（概率循环、小上下文窗口）也预示了 LLM 面临的共性挑战。

